\section{The original coherent quotient and the derived classes it does not remove}
\label{sec:original-coherent-comparison}

This section continues the actual balanced realization in A1427 and the
absolute pointed-base comparison in A1523.  The balanced-kernel chapter,
(BK.1)--(BK.24), proves the original analytic input used below, including
the full-plane Euler inverse and all source-label transition maps.
The derived-packet chapter, (DP.1)--(DP.8), retains both degrees of every
finite spectral thickening.  We compute a different observation here:
maps from the same balanced object into coherent analytic modules, and
then derived maps into those modules.  The distinction is established by
explicit morphisms and resolutions.

\subsection{The original sheaf and its canonical local summands}

Retain the original definitions
\[
\begin{split}
 V&=\{\phi\in\mathcal S(\mathbb R):\phi(-x)=\phi(x),\quad
       \phi(0)=0,\quad\textstyle\int_{\mathbb R}\phi(x)\,dx=0\},\\
 \mathscr B&=\{F\in C^\infty(\mathbb R_{>0}):
   \sup_{x>0}x^b|D^kF(x)|<\infty\quad(b\in\mathbb Z,\ k\ge0)\},
       \qquad D=-x\partial_x,\\
 \Theta\phi(x)&=2\sum_{n\ge1}\phi(nx),\qquad
 \mathcal MF(s)=\int_0^\infty F(x)x^s\frac{dx}{x},\qquad
 Q=\mathscr B/\Theta V,\qquad A=\mathbb C[t],\\
 g(s)&=2\xi(s)=s(s-1)\pi^{-s/2}\Gamma(s/2)\zeta(s),\qquad
 H_\phi(s)=\frac{2\pi^{s/2}M_+\phi(s)}{s(s-1)\Gamma(s/2)}.
\end{split}                                                     \tag{OCQ.1}
\]
Here $t$ acts by $D$ on the original function spaces and by multiplication
by $s$ on analytic functions.  The expression defining $H_\phi$ is its
entire continuation, with every removable point treated in (BK.1)--(BK.8).
In particular $\mathcal M\Theta\phi=gH_\phi$ and
$\mathcal M DF=s\mathcal MF$ on the original spaces.
The tensor products and module Hom groups below are algebraic, as in
the displayed balanced construction in A1427.  The original Schwartz
spaces remain the source of their specified elements and maps.

Let $X=\mathbb C$ with its analytic structure.  Define $\mathcal M$ to
be the sheafification of the presheaf
$U\mapsto\mathcal O_X(U)\otimes_AQ$, and put
$\mathcal N=\mathcal O_X/(g)$.  The presheaf maps
\[
 a\otimes[F]\longmapsto [a\mathcal MF]
\]
are balanced, respect restrictions, and vanish on the original theta
boundaries by (OCQ.1).  They therefore define an actual sheaf morphism
\[
 \beta:\mathcal M\longrightarrow\mathcal N,
 \qquad\mathcal K=\ker\beta.                                  \tag{OCQ.2}
\]
Tensor products commute with filtered colimits, so its stalk at $s_0$ is
the original map
\[
 \begin{gathered}
 O=\mathcal O_{X,s_0},\quad z=s-s_0,\quad
 F_{\mathrm{mer}}=\operatorname{Frac}O=O[1/z],\\
 M=O\otimes_AQ\xrightarrow{\ b\ }N=O/(g),\qquad
 b(a\otimes[F])=[a\mathcal MF],\quad K=\ker b.
 \end{gathered}                                               \tag{OCQ.3}
\]
The letter $F_{\mathrm{mer}}$ denotes meromorphic germs and does not
denote a test function or the formal Laurent-series field.
In these fixed coordinates $O=\mathbb C\{z\}$.

For completeness, surjectivity is visible on an original global test:
\[
 F_0(x)=e^{-(\log x)^2},\qquad
 H_0(s)=\mathcal MF_0(s)=\sqrt\pi e^{s^2/4},\qquad q_0=[F_0].
                                                                    \tag{OCQ.4}
\]
Euler derivatives and multiplication by $x^b$ preserve the required
Gaussian bounds, so $F_0\in\mathscr B$.  Substitution $x=e^y$ evaluates
the Mellin integral, first for real $s$ by completing the square and then
everywhere by entire continuation.  Since $H_0$ is nowhere zero,
$a\mapsto aH_0^{-1}\otimes q_0$ lifts any local section $[a]$ of
$\mathcal N$.  Thus (OCQ.2) is a sheaf epimorphism.

The proved original comparison (BK.14)--(BK.19) says that every nonzero
polynomial acts bijectively on $K$.  Since every nonzero element of $O$
is a unit times a power of $z$, it follows that $K$ is a vector space
over $F_{\mathrm{mer}}$, with the uniquely forced action
$(a/c)k=a(c|_K)^{-1}k$.
It is nonzero at every stalk.  The explicit witness is
\[
 k_0=g\otimes q_0\ne0\quad\text{in }K.                         \tag{OCQ.5}
\]
Here the proof of nonvanishing is exactly the proof retained in
(BK.21): the classical infinitude of distinct actual zeros of $g$
forces $P=0$ in any identity $P(s)H_0(s)=g(s)H_\phi(s)$; hence
$Aq_0\simeq A$, and flatness of the torsion-free $A$-module $O$ gives
$Oq_0\simeq O$.  The classical zero-set input is the retained Hardy
theorem and its original source identified in that chapter.  Neither
RH nor an off-line zero is assumed.  The structure and factorization
arguments below only require the computed vector-space property;
the assertions of noncoherence and nonzero derived classes also use
this proved witness.

Retain the full local factor
$g(s)=u(z)z^m$ with $u(0)\ne0$ and $m=\operatorname{ord}_{s_0}g$.
There is an explicit canonical decomposition.  For $x\in M$, $gx\in K$,
so define
\[
 p_K(x)=(g|_K)^{-1}(gx),\qquad p_T(x)=x-p_K(x),\qquad
 T=M[g].                                                       \tag{OCQ.6}
\]
These are $O$-linear, $p_K$ is the identity on $K$, and $gp_T(x)=0$.
The intersection $T\cap K$ is zero because $g$ is invertible on $K$.
The map $b|_T$ is injective and is surjective because
$b(p_Tx)=b(x)$.  Consequently, with $j=(b|_T)^{-1}$,
\[
 M=T\oplus K,\qquad T\xrightarrow[\ b\ ]{\sim}O/(g),\qquad
 x=j(bx)+p_Kx.                                                 \tag{OCQ.7}
\]
An $O$-torsion element of $M$ has zero $K$ component, so $T$ is the
entire torsion submodule.  This includes $m=0$, when $N=T=0$.
Formula (OCQ.6) retains $u$; no replacement of the function $g$ has
been made.

\subsection{Every ordinary coherent observation factors through the actual zero scheme}

Call an $O$-module divisible when multiplication by every nonzero
element of $O$ is surjective, and call a module reduced when it has
no nonzero divisible submodule.  Over this ring divisibility is
equivalent to surjectivity of multiplication by $z$.
Every finitely generated $O$-module $E$ is reduced.  To see this without
an assumption on a submodule, the elementary divisor decomposition gives
$E\simeq O^r\oplus\bigoplus_{j=1}^d O/(z^{n_j})$.
This decomposition follows by reducing a finite presentation matrix
over the discrete valuation ring: choose a nonzero entry of least
valuation, use its divisibility into all entries to clear its row and
column by invertible operations, and repeat on the remaining finite
matrix.  The resulting diagonal entries are units times powers of $z$;
unit row operations remove the units.  The presentation exists because
every nonzero ideal of $O$ contains an element of smallest valuation
and is generated by it.  More explicitly, for a submodule of $O^n$,
its projection to the first coordinate is a principal ideal.  Choose
one lift of its generator; subtracting its multiples leaves a submodule
of $O^{n-1}$.  Induction on $n$ proves finite generation of every such
submodule and hence of the kernel of a finite free surjection.
The decomposition proves
\[
 \bigcap_{n\ge0}z^nE=0,                                      \tag{OCQ.8}
\]
since a holomorphic germ divisible by every power is zero, and sufficiently
large powers kill each torsion summand.  Every divisible submodule is
contained in this intersection.

An $O$-linear image of $K$ is divisible.  Therefore every map $M\to E$
with $E$ reduced vanishes on $K$ and factors uniquely through $b$.
For a finite module this yields, naturally in $E$,
\[
 b^*:\operatorname{Hom}_O(N,E)\xrightarrow{\sim}
       \operatorname{Hom}_O(M,E),\qquad
 \operatorname{Hom}_O(M,E)\simeq E[g].                         \tag{OCQ.9}
\]
The last map sends $f$ to $f(j([1]))$.  Its inverse sends
$e\in E[g]$ to $x\mapsto a e$ whenever $b(x)=[a]$;
changing $a$ by $gc$ does not change the answer.

The reduced targets form the exact largest class with this property.
Indeed, if an arbitrary target $E$ contains a nonzero divisible submodule
$D$, choose $d_0\ne0$ and $d_{n+1}\in D$ with $zd_{n+1}=d_n$.
The formula
\[
 \alpha:F_{\mathrm{mer}}\longrightarrow D,\qquad
          \alpha(az^{-n})=ad_n                               \tag{OCQ.10}
\]
is well-defined: equality $az^{-n}=bz^{-m}$, with $m\ge n$,
means $b=z^{m-n}a$, while $z^{m-n}d_m=d_n$.
It is additive, $O$-linear, and nonzero at $1$.
Choose an $F_{\mathrm{mer}}$-basis of $K$ containing $k_0$, and let
$\ell:K\to F_{\mathrm{mer}}$ be its coordinate projection with
$\ell(k_0)=1$.  Then $\alpha\ell p_K:M\to E$ is nonzero on $K$,
so it does not factor through $b$.  The choices are explicit choices
of a divisible chain and a vector-space complement, not canonical maps.

Also $K$ is the maximal divisible submodule of $M$: its image under $b$
in the finite module $N$ must be zero.  The universal property implies
\[
 K=\bigcap_{\substack{E\text{ finite over }O\\f:M\to E}}
       \ker f.                                                \tag{OCQ.11}
\]
The inclusion from left to right is (OCQ.9), and the reverse inclusion
uses $f=b$ itself.

Neither $K$ nor $M$ is finitely generated at any analytic stalk.
Otherwise (OCQ.8), applied respectively to $K$ or $M$, contradicts
$0\ne k_0\in\bigcap_nz^nK\subset\bigcap_nz^nM$.
Hence $\mathcal K$ and $\mathcal M$ are nowhere coherent, although
$\mathcal N$ is coherent.  For a coherent analytic sheaf $\mathcal E$,
any actual sheaf morphism $\mathcal M|_U\to\mathcal E|_U$ vanishes
on $\mathcal K|_U$ by (OCQ.9) on stalks.  The sheaf quotient property
gives a unique factorization on $U$, compatible on overlapping opens.
We obtain a proved sheaf isomorphism
\[
 \mathcal H om_{\mathcal O_X}(\mathcal N,\mathcal E)
   \xrightarrow{\sim}
 \mathcal H om_{\mathcal O_X}(\mathcal M,\mathcal E).           \tag{OCQ.12}
\]
Thus $\beta$ is the universal coherent quotient of this original object.
This statement neither identifies $\mathcal M$ with its quotient nor
asserts existence of such a quotient for arbitrary analytic sheaves.

Every coherent quotient of $M$ is consequently a quotient of $O/(g)$.
At the fixed stalk its kernel lifts to an ideal $I$ of $O$ containing
$(g)=(z^m)$.  The valuation argument above gives $I=(z^r)$ for exactly
one integer $0\le r\le m$.  Conversely every such ideal defines a
quotient.  The full map and its removed part are
\[
 M\xrightarrow b O/(g)\longrightarrow O/(z^r),\qquad
 \ker\bigl(O/(g)\to O/(z^r)\bigr)=(z^r)/(g).                  \tag{OCQ.13}
\]
This records all $m-r$ removed nilpotent levels and retains the actual
ideal $(g)$.  Globally these kernels glue exactly to ideal sheaves
between $(g)$ and $\mathcal O_X$, or equivalently to effective
subdivisors of the actual divisor of $g$.
In particular the ordinary analytic duals are
\[
 \operatorname{Hom}_O(M,O)=0,\qquad
 \operatorname{Hom}_O(K,O)=0,\qquad M^{\vee\vee}=0.            \tag{OCQ.14}
\]
The next calculation constructs the derived comparison that (OCQ.14)
does not address.

\subsection{A complete local derived calculation}

All Ext groups in this subsection are groups of modules over the specified
analytic stalk $O=\mathbb C\{z\}$.  Put
$\widehat E=\varprojlim_nE/z^nE$ for a finite $O$-module $E$.
Equation (OCQ.8) embeds $E$ into its completion.
The original source is not replaced by a completed module.

The meromorphic germ field has the following explicit two-term free
resolution.  With $P_0=P_1=\bigoplus_{n\ge0}Oe_n$, set
\[
 0\longrightarrow P_1\xrightarrow d P_0
   \xrightarrow\epsilon F_{\mathrm{mer}}\longrightarrow0,
 \qquad d(e_n)=e_n-ze_{n+1},\quad\epsilon(e_n)=z^{-n}.          \tag{OCQ.15}
\]
Surjectivity follows from the definition of $O[1/z]$.
If a finitely supported vector has last nonzero coefficient $a_m$,
the coefficient of $e_{m+1}$ in its image under $d$ is $-za_m\ne0$;
hence $d$ is injective.  Modulo $dP_1$ one has
$e_n=z^{m-n}e_m$ for $n\le m$.  Therefore a vector
$\sum_{n=0}^m a_ne_n$ in $\ker\epsilon$ is congruent to
$(\sum_{n=0}^m z^{m-n}a_n)e_m=0$.  This proves exactness with the
displayed maps.

Applying $\operatorname{Hom}_O(-,E)$ gives the product map
\[
 \Delta:\prod_{n\ge0}E\longrightarrow\prod_{n\ge0}E,
 \qquad\Delta(a)_n=a_n-za_{n+1},\qquad
 \operatorname{Ext}^1_O(F_{\mathrm{mer}},E)=\operatorname{coker}\Delta.
                                                               \tag{OCQ.16}
\]
The kernel is $\operatorname{Hom}_O(F_{\mathrm{mer}},E)=0$ by
(OCQ.8), and all higher Ext groups vanish by this free resolution.
Define the convergent-in-the-completion sum
\[
 S:\prod_{n\ge0}E\longrightarrow\widehat E/E,
       \qquad S(c)=\left[\sum_{n\ge0}z^nc_n\right].           \tag{OCQ.17}
\]
Every summand of index at least $n$ lies in $z^n\widehat E$, so the
sum is well-defined.  Finite partial sums for $c=\Delta a$ are
$a_0-z^{N+1}a_{N+1}$; the second term tends to zero in the completion.
Thus $S\Delta=0$.

For $E=O$, suppose $S(c)=0$ and write
$a=\sum_{n\ge0}z^nc_n\in O$.  Define
\[
 a_n=z^{-n}\left(a-\sum_{j<n}z^jc_j\right).                  \tag{OCQ.18}
\]
The parenthesis is an analytic germ whose formal Taylor coefficients
below degree $n$ vanish.  Its analytic zero order is consequently at
least $n$, so $a_n\in O$.  Subtraction gives
$a_n-za_{n+1}=c_n$, proving $\ker S=\operatorname{im}\Delta$.
Given an arbitrary formal series $\sum b_nz^n$, choose the constant
germs $c_n=b_n$; this proves surjectivity of $S$.  The argument applies
componentwise to finite free modules.  For a finite torsion module
$T$ killed by $z^q$, one has $\widehat T=T$ and the explicit formula
\[
 a_n=\sum_{j=0}^{q-1}z^jc_{n+j},\qquad
 a_n-za_{n+1}=c_n-z^qc_{n+q}=c_n                             \tag{OCQ.19}
\]
proves surjectivity of $\Delta$.  The finite-module decomposition used
in (OCQ.8) completes the proof for every finite $E$.  Since the map
$S$ itself commutes with every homomorphism $E\to E'$, the resulting
isomorphism is natural, independently of the decomposition used to
verify it:
\[
 \operatorname{Ext}^1_O(F_{\mathrm{mer}},E)
    \xrightarrow[\ S\ ]{\sim}\widehat E/E.                  \tag{OCQ.20}
\]

Choose an $F_{\mathrm{mer}}$-basis $B$ of $K$.  The direct sum over $B$
of (OCQ.15) is a free resolution of $K$.  Applying Hom gives products
of (OCQ.16).  Products of these surjective image-to-cokernel maps are
surjective, and their kernels are the corresponding products of images;
thus their cokernel is the product of the individual cokernels.
Consequently
\[
 \operatorname{Hom}_O(K,E)=0,\qquad
 \operatorname{Ext}^1_O(K,E)\simeq
       \prod_{b\in B}(\widehat E/E),\qquad
 \operatorname{Ext}^i_O(K,E)=0\quad(i\ge2).                   \tag{OCQ.21}
\]
Only the product identification depends on the chosen basis.
The canonical summands in (OCQ.7), together with the resolution
$0\to O\xrightarrow{g}O\to N\to0$, now give
\[
\begin{split}
 \operatorname{Hom}_O(M,E)&\simeq E[g],\\
 \operatorname{Ext}^1_O(M,E)&\simeq
       E/gE\ \oplus\ \operatorname{Ext}^1_O(K,E),\\
 \operatorname{Ext}^i_O(M,E)&=0\qquad(i\ge2).
\end{split}                                                   \tag{OCQ.22}
\]
The first summand in degree one is exactly the map $b^*$ from
$\operatorname{Ext}^1_O(N,E)$.

These statements determine the exact scope for all bounded coherent
derived targets.  If $C\in D^b_{\mathrm{fg}}(O)$, then
\[
 \begin{gathered}
 R\operatorname{Hom}_O(N,C)\xrightarrow{\ b^*\ }
 R\operatorname{Hom}_O(M,C)\text{ is a quasi-isomorphism}\\
 \Longleftrightarrow\quad
 H^q(C)\text{ has finite length for every }q.
 \end{gathered}                                               \tag{OCQ.23}
\]
Indeed, the two-term projective resolution of $K$ and the finite
cohomological filtration of $C$ give the hyper-Ext spectral sequence
with terms $\operatorname{Ext}^p_O(K,H^q(C))$.
The column $p=0$ is zero by (OCQ.21), columns $p\ge2$ are zero by
the resolution, and no differential enters or leaves the remaining
column.  The finite filtration therefore has precisely one possible
nonzero graded piece in each total degree, giving
\[
 H^nR\operatorname{Hom}_O(K,C)
   \simeq\operatorname{Ext}^1_O(K,H^{n-1}(C)).                 \tag{OCQ.24}
\]
For a finite module $E$ the quotient $\widehat E/E$ is zero exactly
when its free rank is zero: its torsion part contributes zero, and
$\mathbb C[[z]]/\mathbb C\{z\}\ne0$, as witnessed by the divergent
formal series $\sum_{n\ge0}n!z^n$.  The ratio of successive absolute
terms at any fixed $z\ne0$ is $(n+1)|z|$, so that series cannot
converge in any nonzero disc.  Since $B$ is nonempty by (OCQ.5),
(OCQ.21), (OCQ.24), and the splitting of $M$ prove both directions
of (OCQ.23).

There is also a concrete nonsplit extension, rather than only an abstract
dimension assertion.  Put $\omega=O\,ds$ and use (OCQ.15).  Define
$\varphi:P_1\to\omega$ by $\varphi(e_n)=n!\,ds$, and set
\[
 E_!=\frac{\omega\oplus P_0}
              {\{(-\varphi(v),d(v)):v\in P_1\}}.
 \qquad
 0\longrightarrow\omega\longrightarrow E_!
       \longrightarrow F_{\mathrm{mer}}\longrightarrow0.      \tag{OCQ.25}
\]
The last map sends $(w,p)$ to $\epsilon(p)$.
Injectivity of the first map follows from injectivity of $d$;
if $\epsilon(p)=0$, write $p=dv$, so the class of $(w,p)$ is the
class of $(w+\varphi(v),0)$.  This proves exactness.
Writing $\bar e_n$ for the class of $e_n$, the relations are
$\bar e_n-z\bar e_{n+1}=n!\,ds$.
A splitting would give compatible lifts
$\bar e_n-a_n\,ds$ and hence
$a_n-za_{n+1}=n!$.  Equations (OCQ.17)--(OCQ.20) would then make
$\sum n!z^n$ convergent, a contradiction.  Thus (OCQ.25) is nonsplit.
Pulling it back along $\ell p_K:M\to F_{\mathrm{mer}}$ from
(OCQ.10) gives a class in $\operatorname{Ext}^1_O(M,\omega)$ whose
restriction to the embedded line $F_{\mathrm{mer}}k_0$ is (OCQ.25).
It is therefore nonzero on $K$ and does not come from
$\operatorname{Ext}^1_O(N,\omega)$.

Accordingly the universal ordinary coherent quotient (OCQ.12) does not
extend to all derived coherent targets.  In particular $\omega[1]$
detects the explicit additional class.  This coexists with the exact
derived tensor comparison (DP.2): for every nonzero germ $h$,
$K\otimes_O^{\mathbf L}O/(h)$ is the contractible two-term complex
$[K\xrightarrow{h}K]$.  The two calculations concern explicitly
different functors.  No identification of stalks of sheaf Ext with
Ext groups of these noncoherent stalks is asserted or needed.

\subsection{The residue morphism into the original pointed-base dual}

The finite analytic summand has a direct relation to the actual A1523
duality.  The resolution using multiplication by the actual $g$ gives
\[
 \operatorname{Ext}^1_O(N,\omega)
       \simeq\omega/g\omega
       \xrightarrow{\sim}g^{-1}\omega/\omega,\qquad
       [v\,ds]\longmapsto[v\,ds/g].                          \tag{OCQ.26}
\]
The second map is well-defined and invertible by multiplication by $g$.
There is a canonical complex-linear residue map
\[
 g^{-1}\omega/\omega\longrightarrow\operatorname{Hom}_{\mathbb C}(N,\mathbb C),
 \qquad [\eta]\longmapsto\bigl([v]\mapsto\operatorname{Res}_{s=s_0}(v\eta)\bigr).
                                                               \tag{OCQ.27}
\]
Adding a holomorphic form or changing $v$ by a multiple of $g$ changes
the integrand by a holomorphic form, so the map is well-defined.
It is perfect: with $g=u(z)z^m$, the basis $1,z,\ldots,z^{m-1}$
is dual to the classes
$u(z)z^{m-1-j}\,ds/g=z^{-j-1}ds$, for $0\le j<m$.
Their residue matrix is the identity.  This calculation retains the
entire analytic unit $u$ in the representative before cancellation.
For $m=0$ both sides are zero.

Let $Z$ be an actual finite set of zeros stable under
$\rho\mapsto\rho^\dagger=1-\overline\rho$, and retain all multiplicities
$m_\rho=\operatorname{ord}_\rho g$ in
$A_Z=\bigoplus_{\rho\in Z}\mathcal O_{X,\rho}/(g)$.
The Mellin map from the original quotient is
$J_Z:Q\to A_Z$, $J_Z[F]=([\mathcal MF]_\rho)_\rho$.
Its surjectivity is the original finite Mellin-jet statement retained
in A1523, with complete interpolation proof in (BK.9)--(BK.11);
no arbitrary surrogate packet is used.
For $f\in A_Z$, use the reflected germ
$f^\dagger(s)=\overline{f(1-\overline s)}$ at the corresponding paired
centre.  The original identities $g(1-s)=g(s)$ and
$\overline{g(\overline s)}=g(s)$ give $g^\dagger=g$, so this operation
preserves the actual ideals and is well-defined on the full packets.
Equations (OCQ.26)--(OCQ.27) give the exact composite
\[
 \overline{A_Z}\longrightarrow
 \bigoplus_{\rho\in Z}\operatorname{Ext}^1_{\mathcal O_{X,\rho}}
           (\mathcal O_{X,\rho}/(g),\mathcal O_{X,\rho}\,ds)
 \longrightarrow\operatorname{Hom}_{\mathbb C}(Q,\mathbb C),
                                                               \tag{OCQ.28}
\]
where its value at $f$ is
\[
 [F]\longmapsto
       \sum_{\rho\in Z}\operatorname{Res}_{s=\rho}
              \frac{f^\dagger(s)\mathcal MF(s)}{g(s)}\,ds.      \tag{OCQ.29}
\]
The first map sends $f$ to $[f^\dagger ds/g]$; it is a complex-linear
isomorphism from the conjugate vector space.  The second map is
injective by the perfectness in (OCQ.27) and surjectivity of $J_Z$.
Thus (OCQ.28) is precisely the finite residue injection retained in
A1523, after equipping its target with the original one-dimensional
dilation character $\mathcal L_1$.

To check that character and every sign, let
$R_aF(x)=F(x/a)$ for $a>0$.  Direct Mellin substitution gives
$\mathcal M R_aF(s)=a^s\mathcal MF(s)$.
The original dual action on $\operatorname{Hom}_{\mathbb C}(Q,\mathcal L_1)$
is $(U_a\lambda)(q)=a\lambda(R_{1/a}q)$.  In (OCQ.29) its effect
is multiplication of the differential representative by $a^{1-s}$.
On the source packet, $(a^sf)^\dagger=a^{1-s}f^\dagger$.
Thus the displayed injection is equivariant, with the exact character
$a$ and the original reflection.
Inserting $g'$ gives the original trace form, since
$g'/g=m_\rho/(s-\rho)+u'/u$ at each centre:
\[
 \sum_{\rho\in Z}\operatorname{Res}_{s=\rho}
           \frac{f^\dagger(s)g'(s)v(s)}{g(s)}\,ds
   =\sum_{\rho\in Z}m_\rho
             \overline{f(1-\overline\rho)}\,v(\rho).          \tag{OCQ.30}
\]
This proves the exact morphism from the finite coherent dual summand
to the actual pointed-base comparison; the additional classes in
(OCQ.25) do not replace or delete that summand.

The A1523 target is the complex-linear right adjoint for the original
four-point amplitude diagram.  Explicitly its points satisfy
$+<\eta<\sigma$ and $-<\eta<\sigma$; the amplitudes are
$\mathcal T_+=\mathcal T_-=V$, $\mathcal T_\eta=\mathscr B$,
$\mathcal T_\sigma=0$, with restrictions $\Theta$ and
$J\Theta=\Theta\widehat{\phantom\phi}$, where
$JF(x)=x^{-1}F(1/x)$ and the Fourier kernel is $e^{-2\pi ix\nu}$.
Its cohomology complex is
$[V\oplus V\xrightarrow{(\phi,\psi)\mapsto\Theta(\phi-\widehat\psi)}
\mathscr B]$, in degrees zero and one, so $H^1=Q$.
The right adjoint computed there satisfies
$H^{-1}R\operatorname{Hom}(\mathcal T,K_\tau(\mathcal L_1))
\simeq\operatorname{Hom}_{\mathbb C}(Q,\mathcal L_1)$.
Consequently ordinary $O$-dual vanishing in (OCQ.14) cannot be inserted
as a vanishing statement about that different, explicitly related
right-adjoint target.

\subsection{Proper labels and the precise consequence for purity comparison}

For every original $D$-invariant source $W\subset V$, keep
\[
 Q_W=\mathscr B/\Theta W,\quad M_W=O\otimes_AQ_W,\quad
 L_W=O\otimes_A(V/W),\quad
 \pi_W:M_W\to M,\quad b_W=b\pi_W.                            \tag{OCQ.31}
\]
The inclusion $L_W\to M_W$ sends $a\otimes[\phi]$ to
$a\otimes[\Theta\phi]$.  Injectivity of the original $\Theta$ and
flatness of $O$ over $A$ prove
\[
 0\to L_W\to M_W\xrightarrow{\pi_W}M\to0,\qquad
 0\to L_W\to\ker b_W\to K\to0.                              \tag{OCQ.32}
\]
For a finite target $E$, a map $f:M_W\to E$ factors through $b_W$
if and only if $f|_{L_W}=0$.  The reverse direction is proved by first
descending along the first quotient in (OCQ.32) and then applying
(OCQ.9); the forward direction follows because $b_WL_W=0$.
More precisely,
\[
 0\to\operatorname{Hom}_O(N,E)\xrightarrow{\ b_W^*\ }
 \operatorname{Hom}_O(M_W,E)\longrightarrow
 \operatorname{Hom}_O(L_W,E)\xrightarrow\delta
 \operatorname{Ext}^1_O(M,E)                                \tag{OCQ.33}
\]
is exact.  Here $\delta(\alpha)$ is the pushout of the first sequence
in (OCQ.32) by $\alpha:L_W\to E$: its middle module is
$(E\oplus M_W)/\{(-\alpha(l),l):l\in L_W\}$.
This extension splits exactly when $\alpha$ extends to $M_W$.
Indeed an extension $f$ defines a retraction by
$[e,x]\mapsto e+f(x)$, and any retraction gives $f(x)=[0,x]$
followed by that retraction, with $f|_{L_W}=\alpha$.
This proves exactness at the last displayed Hom term, while the
previous factorization argument proves exactness at the others.

If $W\subset W'\subset V$, the source-quotient maps give the commuting
diagram of (OCQ.32), with
$L_W\to L_{W'}$ induced by $V/W\to V/W'$ and
$M_W\to M_{W'}$ induced by $Q_W\to Q_{W'}$.
Their respective kernels are both
$O\otimes_A(W'/W)$, identified by $[\phi]\mapsto[\Theta\phi]$.
This follows by tensoring the two exact original quotient sequences.
Thus the coherent observations of a proper source retain exactly their
extendable observations on $L_W$ and the obstruction map $\delta$;
they are not replaced by the universal property for the full source.
For example, the actual proper-label classes constructed in
(PL.8)--(PL.15) lie in $L_W$ and can have finite spectral torsion that
vanishes only under this specified transition.

All these maps act inside the original support fibres.  For any fixed
nonempty support label $S$, the split-zero extension of a linear map
$f:E\to E'$ is
\[
 \tau\longmapsto\tau,\qquad
 (S,v)\longmapsto(S,f(v)),\qquad e_S=(S,0).                   \tag{OCQ.34}
\]
Its inverse image of $e_S$ is the full supported kernel
$\{(S,v):v\in\ker f\}$, and its inverse image of $\tau$ is only
$\{\tau\}$.  This follows directly from the disjoint union defining
the split-zero object.  Formula (OCQ.34) applies to $\beta$, to each
label transition, and to every quotient in (OCQ.13); no nonempty
kernel fibre is identified with external absence.

At an actual zero $\rho$ of order $m$, scalar multiplication by
$a^s$ on $N$ is
$a^\rho\exp((\log a)z)$ modulo the full ideal $(g)$.
For the quotient of depth $r$ in (OCQ.13), the removed module
$(z^r)/(g)$ has the ordered basis
$z^r,\ldots,z^{m-1}$, and the action is triangular with diagonal
$a^\rho$.  Therefore its dimension and trace are exactly
\[
 \dim_{\mathbb C}(z^r)/(g)=m-r,
 \qquad\operatorname{Tr}\bigl(a^s;(z^r)/(g)\bigr)
                 =(m-r)a^\rho.                              \tag{OCQ.35}
\]
Every partial or total deletion of an actual local packet is consequently
recorded by a nonzero finite coherent kernel and by this exact change in
the original trace, including all multiplicities.
The dilation on the full balanced module remains the original
$1\otimes R_a$, whose difference from scalar $a^s$ has image in $K$
as proved in (DP.4)--(DP.6).  Formula (OCQ.35) is asserted on the actual
analytic packet and its specified quotient, where the actions agree.

The obtained consequence for the tau purity comparison is exact.
The full original quotient has the actual zero scheme as its universal
ordinary coherent quotient; retaining that universal map retains every
actual off-line packet if such a packet exists.  Proper source labels
retain the additional maps and kernels in (OCQ.32)--(OCQ.33).
Derived coherent targets of positive rank additionally see the explicit
kernel extension (OCQ.25), whereas every finite-length derived target
has exactly the comparison (OCQ.23).  The original finite residue
injection is the composite (OCQ.28), with its computed action and trace.
None of these proved maps supplies a positive metric or a weight bound
on the actual packet; no numerical location constraint has been deduced
from coherence.  The original purity question therefore continues with
these full morphisms and obstructions retained, rather than with a
claimed identification of the full analytic object and its coherent
quotient.
